\chapter{Acromegaly}
\index{Acromegaly}
\label{sec:Acromegaly}

\section{Overview}
Acromegaly is a chronic disorder caused by excessive growth hormone (GH) secretion, almost always from a pituitary somatotroph adenoma (95\%), with GH stimulating liver production of insulin-like growth factor 1 (IGF-1) which mediates most of the somatic effects, and onset being insidious with a mean delay of 7--10 years between symptom onset and diagnosis.
\section{Causes}
The underlying mechanisms involve complex interactions between genetic predisposition and environmental factors. Disruption of normal physiological processes leads to the characteristic features of the condition. Multiple contributing factors may act synergistically to trigger or worsen the condition. Understanding the specific cause helps guide targeted treatment approaches.
\section{Risk Factors}


\begin{itemize}[noitemsep]
  \item Genetic predisposition and family history
  \item Advancing age
  \item Lifestyle factors (diet, exercise, smoking, alcohol)
  \item Environmental exposures
  \item Underlying medical conditions
  \item Medication use
\end{itemize}
\section{Signs and Symptoms}

\subsection{Common symptoms}
\begin{itemize}[noitemsep]
  \item You may experience coarsening of facial features, prognathism (jaw protrusion), frontal bossing, enlarged hands and feet, soft tissue swelling, macroglossia, excessive sweating, deepening of the voice, arthropathy, carpal tunnel syndrome, high blood pressure, diabetes mellitus, sleep apnea, cardiomyopathy, and increased risk of colorectal polyps and cancer. Diagnosis shows elevated age- and sex-matched IGF-1 level and failure to suppress GH below 1 ng/mL after 75-g oral glucose tolerance test, with MRI showing a pituitary macroadenoma in most cases. Symptoms vary depending on the severity and extent of the condition Pain or discomfort in the affected area Functional impairment affecting daily activities Changes in appearance or function of affected tissues Systemic symptoms may include fatigue, fever, or weight changes Symptoms may develop gradually or appear suddenly
  \item Pain or discomfort in the affected area
  \end{itemize}

\subsection{Emergency warning signs}
\begin{itemize}[noitemsep]
  \item Severe or worsening symptoms of acromegaly require immediate medical evaluation
\end{itemize}
\section{Progression and Stages}
The condition may progress through different stages of severity over time. Early stages often involve mild or subtle symptoms that may be overlooked. Moderate stages involve more noticeable symptoms affecting daily function. Advanced stages may involve significant impairment and complications.
\section{Complications}
\begin{itemize}[noitemsep]
  \item Disease progression leading to worsening symptoms
  \item Secondary infections or injuries
  \item Side effects of treatment
  \item Reduced quality of life
  \item Emotional and psychological impact
\end{itemize}
\section{Diagnosis}
Comprehensive medical history and physical examination Laboratory tests to assess organ function and rule out other conditions Imaging studies to visualize affected structures Specialized testing depending on the affected system Consultation with relevant specialists Monitoring of disease progression over time

\begin{itemize}[noitemsep]
  \item Comprehensive medical history and physical examination
  \item Laboratory tests to assess organ function
  \item Imaging studies to visualize affected structures
  \item Specialized testing depending on the affected system
  \item Consultation with relevant specialists
\end{itemize}
\section{Prevention}
\begin{itemize}[noitemsep]
  \item Maintain a healthy lifestyle with balanced diet and exercise
  \item Avoid known risk factors
  \item Attend regular medical check-ups for early detection
  \item Follow recommended screening guidelines
\end{itemize}
\section{Treatment Options}
Treatment options include transsphenoidal surgery (first-line, curative in 80--90\% of microadenomas), somatostatin analogs (octreotide LAR, lanreotide, pasireotide) for persistent disease, pegvisomant (GH receptor antagonist), and dopamine agonists (cabergoline) for mild cases. Medications to manage symptoms and address underlying mechanisms Lifestyle modifications and self-care measures Physical therapy and rehabilitation Surgical intervention when indicated Regular monitoring and follow-up care Patient education and support services

\begin{itemize}[noitemsep]
  \item Medications to manage symptoms and address underlying mechanisms
  \item Lifestyle modifications and self-care measures
  \item Physical therapy and rehabilitation
  \item Surgical intervention when indicated
  \item Regular monitoring and follow-up care
\end{itemize}
\section{Living with It}
\begin{itemize}[noitemsep]
  \item Follow your treatment plan as prescribed
  \item Maintain regular follow-up with your healthcare provider
  \item Make necessary lifestyle adjustments
  \item Seek support from family, friends, or support groups
  \item Stay informed about your condition
\end{itemize}
\section{Outlook}
outlook is improved with early treatment, which reduces cardiovascular morbidity and mortality. The outlook varies depending on the specific condition, severity at diagnosis, and response to treatment. Early intervention generally leads to better outcomes. Many conditions can be effectively managed with appropriate treatment. Regular follow-up care is important for monitoring and treatment adjustment.
